\section{Peripheral automorphisms and the pronilpotent quotient}
\label{sec:peripheral-lifting}

We prove a lifting theorem with prescribed peripheral conjugators.
Keeping these conjugators is necessary when a peripheral element must
remain fixed after lifting. The finite argument uses central extensions
to retain the character values which control the relevant lifting fibres.

There are two source groups. In the free case set
\[
 \Lambda=\widehat F_2=\langle x_1,x_2\rangle,
 \qquad x_3=(x_1x_2)^{-1},\qquad U=\Lambda.
\]
In the modular case set
\[
 \Lambda=\widehat{C_2*C_3}
     =\langle x_1,x_2\mid x_1^2=x_2^3=1\rangle,
 \qquad x_3=(x_1x_2)^{-1},
\]
and let $U$ be the kernel of the standard quotient $\Lambda\to S_3$.
The theorem below also holds for any closed normal subgroup $U$ of
either source; these two choices are its initial applications.
For such a subgroup, let $\gamma_n(U)$ be its closed lower central series,
and put
\[
 R=\bigcap_{n\geq1}\gamma_n(U),\qquad
 \overline\Lambda=\Lambda/R,
 \qquad \pi:\Lambda\longrightarrow\overline\Lambda.
\]
A positive peripheral automorphism sends each $x_i$ to a conjugate of
$x_i$, with no change of exponent.

\begin{theorem}[Peripheral lifting]\label{thm:peripheral-lifting}
Let $\alpha$ be a continuous automorphism of $\overline\Lambda$.
Suppose that elements $\overline c_i\in\overline\Lambda$ satisfy
\[
 \alpha(\pi(x_i))
    =\overline c_i\pi(x_i)\overline c_i^{-1}
       \qquad(1\leq i\leq3).
\]
There are a continuous automorphism $\Phi$ of $\Lambda$ and elements
$c_i\in\Lambda$ such that
\[
 \pi\Phi=\alpha\pi,
 \qquad \pi(c_i)=\overline c_i,
 \qquad \Phi(x_i)=c_ix_ic_i^{-1}
       \qquad(1\leq i\leq3).
\]
If one specified conjugator $\overline c_i$ is the identity, the lift
can be chosen so that the corresponding $c_i$ is the identity.
\end{theorem}

The last assertion normalizes one chosen conjugator. It does not claim
that several full conjugators can be normalized independently by a
single inner automorphism.

\subsection{A finite central refinement}

\begin{lemma}\label{lem:finite-peripheral-refinement}
Let $q:\Lambda\twoheadrightarrow E$ be a finite quotient, let
$A\triangleleft E$, and write $D=E/A$ and $q_D=q\bmod A$.
There are a finite central extension
$\beta:D^\dagger\twoheadrightarrow D$ and an epimorphism
$q^\dagger:\Lambda\twoheadrightarrow D^\dagger$ such that
$\beta q^\dagger=q_D$, with the following property.

Suppose $q_1^\dagger:\Lambda\twoheadrightarrow D^\dagger$ is an
epimorphism and $r_i\in D^\dagger$ satisfy
\[
 q_1^\dagger(x_i)=r_iq^\dagger(x_i)r_i^{-1}
       \qquad(1\leq i\leq3).
\]
There is an epimorphism $q_1:\Lambda\twoheadrightarrow E$ and
elements $C_i\in E$ with
\[
 q_1(x_i)=C_iq(x_i)C_i^{-1},\qquad
 C_i\bmod A=\beta(r_i),\qquad
 q_1\bmod A=\beta q_1^\dagger.
\]
\end{lemma}

\begin{proof}
Put $X_i=q(x_i)$ and let
\[
 T_i=\{aX_ia^{-1}:a\in A\}
\]
be its $A$-conjugacy class. We shall retain enough information to
compare homomorphism counts into every subgroup of $E$, and then
recover the number of epimorphisms.

Fix a subgroup $J\leq E$ mapping onto $D$ such that each $T_i\cap J$
is nonempty, and put $M=J\cap A$. Let
$\rho:M\to\operatorname{GL}(V_\rho)$ be an irreducible complex
representation whose isomorphism class is fixed by conjugation by $J$.
Write $d_\rho=\dim V_\rho$ and $e=\exp J$. Define
\[
 \widetilde J_\rho=
 \left\{(j,B)\in J\times\operatorname{GL}(V_\rho):
 \begin{array}{l}
 B\rho(m)B^{-1}=\rho(jmj^{-1})\quad(m\in M),\\
 \det B\in\mu_e
 \end{array}\right\}.
\]
The invariance of $\rho$ gives an implementing matrix above every $j$.
By Schur's lemma all such matrices form one scalar coset. In that
coset the determinant condition has exactly $d_\rho e$ solutions:
if $B_0$ is one implementing matrix, it is the equation
$t^{d_\rho e}=(\det B_0)^{-e}$ for $B=tB_0$. Matrix multiplication
and inversion preserve the defining conditions. Thus
$\widetilde J_\rho$ is a finite group, and its kernel over $J$ is
the central scalar group $\mu_{d_\rho e}$.

The map
\[
 M\longrightarrow\widetilde J_\rho,
 \qquad m\longmapsto(m,\rho(m))
\]
is a normal embedding. Its matrices satisfy the determinant condition
because $m^e=1$, and normality is exactly the intertwining identity.
This copy of $M$ intersects the scalar kernel trivially. Consequently
\[
 Q_\rho=\widetilde J_\rho/M
\]
is a finite central extension of $J/M=D$. The natural map gives an
isomorphism
\begin{equation}
 \widetilde J_\rho\simeq J\times_DQ_\rho.
 \label{eq:representation-fibre-product}
\end{equation}
Indeed it is injective, and any pair with matching images in $D$ is
obtained by multiplying a representative in $\widetilde J_\rho$ by
an element of the embedded $M$. The matrix coordinate therefore gives
a representation
\[
 \Psi_\rho:J\times_DQ_\rho
        \longrightarrow\operatorname{GL}(V_\rho)
\]
whose restriction to $M\times\{1\}$ is $\rho$.

We next choose a source homomorphism
$R_\rho^0:\Lambda\to Q_\rho$ lifting $q_D$. For the free source,
choose arbitrary lifts of $q_D(x_1)$ and $q_D(x_2)$. For the modular
source choose $j_i\in J\cap T_i$ for $i=1,2$. The orders of $j_1$
and $j_2$ divide two and three, respectively. If $j_i$ has order
$h>1$, an implementing matrix has its $h$th power scalar by Schur's
lemma. Rescaling makes that power the identity. Its determinant then
has order dividing $h$, and hence lies in $\mu_e$. If $j_i=1$, use
the identity matrix. The resulting elements of $Q_\rho$ have orders
dividing two and three, respectively, and define the required
homomorphism from the modular source.

Perform this construction for all the finitely many relevant $J$ and
all their invariant irreducible representations. Take the joint image
of the maps $q_D$ and $R_\rho^0$ in their fibre product over $D$.
Call it $D^\dagger$, and call the resulting source epimorphism
$q^\dagger$. Its kernel over $D$ is a subgroup of a product of central
scalar groups, so $D^\dagger\to D$ is a finite central extension.

We now establish the counting identity used to transport this data.
Fix one of the subgroups $J$ above. Let $\chi:\Lambda\to D$ be
an epimorphism. Suppose $V_i\subset J$ is invariant under conjugation
by $M$ and lies in the fibre over $\chi(x_i)$. For each invariant
$\rho$ suppose $R_\rho:\Lambda\to Q_\rho$ is a homomorphism lifting
$\chi$. Then the number of triples with product one is
\begin{equation}
 \begin{split}
 &\#\{(a_1,a_2,a_3)\in V_1\times V_2\times V_3:
                   a_1a_2a_3=1\}\\
 &\qquad=\frac1{|M|}
  \sum_{\rho\in\operatorname{Irr}(M)^J}\frac1{d_\rho}
  \prod_{i=1}^3
    \left(\sum_{a\in V_i}
       \operatorname{tr}\Psi_\rho(a,R_\rho(x_i))\right).
 \end{split}
 \label{eq:relative-peripheral-count}
\end{equation}
To prove this, put $B_i=\sum_{a\in V_i}a$ in $\mathbb C[J]$.
Each $B_i$ commutes with $\mathbb C[M]$. Let $e_\tau$ be the
primitive central idempotent of $\mathbb C[M]$ belonging to an
irreducible representation $\tau$. Since every term of $B_i$ has
the same image in $D$, multiplication by $B_i$ takes the $\tau$
block to the block obtained by conjugating $\tau$ by that image.
But $e_\tau B_i=e_\tau B_ie_\tau$, by commutation. Hence this
product is zero unless $\chi(x_i)$ fixes $\tau$. The images of
$x_1,x_2$ generate $D$, so only $J$-invariant representations can
contribute to $B_1B_2B_3$.

This product lies in $\mathbb C[M]$, because the three quotient
values multiply to one. The coefficient of the identity in an
element $B\in\mathbb C[M]$ is
\[
 \frac1{|M|}\sum_{\tau\in\operatorname{Irr}(M)}
                  (\dim\tau)\operatorname{tr}\tau(B),
\]
by the regular representation of $M$. For an invariant $\rho$, set
\[
 S_{i,\rho}=\sum_{a\in V_i}\Psi_\rho(a,R_\rho(x_i)).
\]
Conjugating by $\rho(m)$ permutes this sum, so Schur's lemma makes
it a scalar matrix, namely
$d_\rho^{-1}\operatorname{tr}(S_{i,\rho})I$.
Since $R_\rho(x_1)R_\rho(x_2)R_\rho(x_3)=1$, multiplication in
the fibre product gives
\[
 \rho(B_1B_2B_3)=S_{1,\rho}S_{2,\rho}S_{3,\rho}.
\]
The three scalar denominators and the two dimension factors in the
regular-character calculation give the factor $1/d_\rho$ in
\eqref{eq:relative-peripheral-count}. This proves the identity.
For the modular source, when $V_1,V_2$ are subsets of the prescribed
conjugacy classes, their elements already satisfy the square and cube
relations. In that case the same triple count is the number of source
homomorphisms with these conditions.

Return now to $q_1^\dagger$ and its specified conjugators $r_i$.
Choose $H_i\in E$ mapping to $\beta(r_i)$ and put
$T_i^1=H_iT_iH_i^{-1}$. For every $J\leq E$ mapping onto $D$,
choose $j_i\in J$ with the same image as $H_i$ in $D$.
Writing $H_i=j_ia_i$ with $a_i\in A$ gives the exact equality
\begin{equation}
 T_i^1\cap J=j_i(T_i\cap J)j_i^{-1}.
 \label{eq:relative-class-intersection}
\end{equation}
Here we use that $T_i$ is an entire $A$-conjugacy class; conjugation
by $a_i$ preserves it, and conjugation by $j_i$ preserves $J$.

For a relevant $J$ and $\rho$, let $r_{i,\rho}$ be the projection
of $r_i$ to $Q_\rho$, and let $R_\rho^1$ be the corresponding
projection of $q_1^\dagger$. The pair $(j_i,r_{i,\rho})$ belongs
to $J\times_DQ_\rho$. It conjugates the old pairs
\[
 (a,R_\rho^0(x_i)),\quad a\in T_i\cap J,
\]
to the new pairs
\[
 (j_iaj_i^{-1},R_\rho^1(x_i)),\quad
       j_iaj_i^{-1}\in T_i^1\cap J.
\]
The trace sums in \eqref{eq:relative-peripheral-count} are therefore
equal term by term. The number of homomorphisms $\Lambda\to J$
with reduction $q_D$ and values in $T_i$ equals the number with
reduction $\beta q_1^\dagger$ and values in $T_i^1$.
If some $T_i\cap J$ is empty, both counts are zero by
\eqref{eq:relative-class-intersection}. If $J$ does not map onto $D$,
both counts are zero because both specified quotient maps are onto.
Thus the equality holds for every subgroup $J\leq E$.

Each homomorphism into $J$ has a unique image subgroup. M\"obius
inversion on the finite subgroup lattice, or induction on subgroup
order, recovers the number of epimorphisms to $E$ from these counts.
The two epimorphism counts agree. The old count is positive because
it contains $q$. Hence a new epimorphism $q_1$ exists. For each $i$,
membership $q_1(x_i)\in T_i^1$ gives an element $a_i'\in A$ such
that
\[
 q_1(x_i)=H_ia_i'X_i(a_i')^{-1}H_i^{-1}.
\]
Taking $C_i=H_ia_i'$ gives the prescribed conjugator image in $D$.
The quotient identity follows on the two generators, proving the lemma.
\end{proof}

\subsection{Nilpotence and compactness}

\begin{proof}[Proof of Theorem~\ref{thm:peripheral-lifting}]
First observe that a continuous map from $\overline\Lambda$ to a
finite group has nilpotent image of $U$. Indeed its kernel on $U$
is an open subgroup containing $R=\bigcap_n\gamma_n(U)$.
Compactness of $U$ and the descending closed family $\gamma_n(U)$
imply that this kernel contains some $\gamma_n(U)$. Conversely,
any finite source homomorphism whose image of $U$ is nilpotent kills
$R$ and factors through $\overline\Lambda$.

Let $q:\Lambda\twoheadrightarrow E$ be a finite quotient. In the
finite group $q(U)$ the lower central series stabilizes. Put
\[
 A=\gamma_\infty(q(U))
       =\bigcap_{n\geq1}\gamma_n(q(U)).
\]
The subgroup $A$ is normal in $E$, and $q(U)/A$ is nilpotent.
Apply Lemma~\ref{lem:finite-peripheral-refinement} to obtain
$q^\dagger:\Lambda\twoheadrightarrow D^\dagger$.
The subgroup $q^\dagger(U)$ is a central extension of $q(U)/A$,
so it is nilpotent as well: a central extension of a group of class
$c$ has class at most $c+1$. It follows that $q^\dagger$ factors
as $\overline q^\dagger\pi$. Define
\[
 q_1^\dagger=\overline q^\dagger\alpha\pi,
 \qquad r_i=\overline q^\dagger(\overline c_i).
\]
This is an epimorphism with the required three conjugacies. The
finite lemma supplies an epimorphism $q_1:\Lambda\to E$ and
conjugators $C_i\in E$ whose images modulo $A$ are the prescribed
images of $\overline c_i$.

We use these finite lifts simultaneously. In the compact space
$\Lambda^3$, impose the following closed conditions on $(c_1,c_2,c_3)$:
\begin{enumerate}
 \item the elements $c_ix_ic_i^{-1}$ have product one;
 \item their first two images generate every finite quotient of $\Lambda$;
 \item the image of each $c_i$ in every finite quotient of
       $\overline\Lambda$ agrees with the image of $\overline c_i$.
\end{enumerate}
For the first condition it is enough to impose its images in finite
quotients. Each finite list of conditions is controlled by one finite
quotient $q:\Lambda\to E$: take the joint image of all source
quotients in the list, including the quotients that factor through
$\overline\Lambda$.

For this joint quotient, $A=\gamma_\infty(q(U))$ maps trivially to
each quotient in the finite list of $\overline\Lambda$, by the nilpotence
observation above. Thus the conjugators $C_i\in E$ from the finite
lemma have all the prescribed images. Arbitrary lifts of the
$C_i$ to $\Lambda$ satisfy the finite list: their peripheral images
have product one in $E$, generate $E$, and have the prescribed
conjugator values in each quotient in the finite list. The family of closed
conditions therefore has the finite intersection property.
Compactness gives a triple satisfying every condition.

The source presentation now defines a continuous endomorphism
$\Phi$ taking $x_i$ to $c_ix_ic_i^{-1}$. In the modular case the
square and cube relations hold because these are conjugates of the
original generators; the product condition gives the required third
image in both cases. Its image maps onto every finite quotient of
$\Lambda$ by the generation conditions. The image is compact and
therefore closed, so $\Phi$ is surjective.

Finally a finitely generated profinite group is Hopfian. To recall the
argument, for each positive integer $d$ there are only finitely many
open normal subgroups of index $d$. Inverse image under a surjective
endomorphism injects this finite set into itself, preserving the index,
and is therefore a permutation. Every such subgroup is consequently
the inverse image of an open normal subgroup and contains the kernel
of the endomorphism. Intersecting all open normal subgroups shows that
the kernel is trivial. Hence $\Phi$ is an automorphism.

The prescribed conjugator conditions imply
$\pi\Phi(x_i)=\alpha\pi(x_i)$ for the two generators, and therefore
$\pi\Phi=\alpha\pi$. If a selected $\overline c_i=1$, then
$c_i\in R$. Replace $\Phi$ by its composite with conjugation by
$c_i^{-1}$ and replace every conjugator $c_j$ by $c_i^{-1}c_j$.
This leaves their images in the quotient and the quotient automorphism
unchanged and makes the selected conjugator equal to one.
\end{proof}
